\section*{ВВЕДЕНИЕ}
\addcontentsline{toc}{section}{ВВЕДЕНИЕ}

В условиях растущего давления санкций и ограничений на использование зарубежного программного обеспечения российским компаниям становится все сложнее поддерживать эффективные внутренние коммуникации. Многие популярные корпоративные мессенджеры либо вводят ограничения для пользователей из России, либо полностью прекращают работу на территории страны. В такой ситуации критически важно иметь независимое решение, не зависящее от зарубежных поставщиков и доступное для использования без риска отключения или потери данных.

Закрытые корпоративные мессенджеры отечественной разработки позволяют компаниям обходить санкционные барьеры, обеспечивая стабильную работу внутренних коммуникаций. Они функционируют без зависимости от западных облачных сервисов, не требуют сторонних API и имеют полное управление данными внутри организации. Такие решения позволяют сохранить бизнес-процессы, поддерживать работу команд и исключить риски внезапных отключений.

Ранее корпоративные сети и коммуникационные платформы в России во многом зависели от западных технологий, что делало их уязвимыми перед санкционными угрозами. Развитие отечественного программного обеспечения позволяет создать гибкие и надежные системы, которые легко адаптируются под нужды компании. Современные российские корпоративные мессенджеры интегрируются с внутренними сервисами предприятия, обеспечивают удобное управление пользователями, защиту данных и полную автономность от зарубежных сервисов.

Главной задачей независимого корпоративного мессенджера является предоставление стабильной платформы для общения, управления бизнес-процессами и обмена файлами без ограничений со стороны иностранных разработчиков.

\textit{Цель настоящей работы} – разработка корпоративного мессенджера, полностью независимого от санкций и обеспечивающего бесперебойную внутреннюю коммуникацию. Для достижения поставленных целей необходимо решить следующие задачи:

\begin{enumerate}
	\item Спроектировать архитектуру серверной части корпоративного мессенджера с использованием Python, WebOb и WSGI, обеспечив маршрутизацию запросов и модульную структуру приложения.
	
	\item Реализовать механизм регистрации, авторизации и управления пользовательскими сессиями с привязкой к уникальному идентификатору и криптографическим ключам, обеспечив безопасность передачи сообщений.
	
	\item Разработать интерфейсы (API) для обработки запросов, связанных с отправкой, получением, редактированием и удалением сообщений в групповых и приватных чатах, включая асинхронную проверку обновлений.
	
	\item Организовать работу с файловыми вложениями: загрузка, сохранение, определение MIME-типов, отображение изображений и предоставление возможности скачивания других типов файлов.
	
	\item Внедрить модуль управления группами: создание, переименование, добавление и исключение участников, управление ролями (участник, администратор, владелец), проверка прав доступа и уведомления о событиях.
	
	\item Обеспечить серверную обработку системных сообщений, их отображение в интерфейсе.
	
	\item Разработать механизм полнотекстового поиска по сообщениям с поддержкой фильтрации по типу чата, ID чата и сортировкой по дате или релевантности.
	
	\item Спроектировать обработку ошибок с возвратом корректных кодов HTTP-ответов, логированием исключений и сохранением отказоустойчивости сервиса.
	
	\item Провести модульное и интеграционное тестирование логики серверной части: проверка корректности маршрутов, безопасности доступа, устойчивости к некорректным запросам и проверка взаимодействия с базой данных.
\end{enumerate}


\textit{Структура и объем работы}. Отчет состоит из введения, 4 разделов основной части, заключения, списка использованных источников, 2 приложений. Текст выпускной квалификационной работы равен \formbytotal{lastpage}{страниц}{е}{ам}{ам}.

\textit{Во введении} сформулирована цель работы, поставлены задачи разработки, описана структура работы, приведено краткое содержание каждого изразделов.

\textit{В первом разделе} на стадии описания технической характеристики предметной области приводится сбор информации о возможностях систем управления ресторанного бизнеса, а также производится анализ существующих аналогов.

\textit{Во втором разделе} на стадии технического задания приводятся требования к разрабатываемому продукту.

\textit{В третьем разделе} на стадии технического проектирования представлены проектные решения для системы управления.

\textit{В четвертом разделе} приводится список классов и их методов, использованных при разработке программы, производится тестирование разработанной системы.

\textit{В заключении} излагаются основные результаты работы, полученные в ходе разработки.

В приложении А представлен графический материал. В приложении Б представлены фрагменты исходного кода.